\documentclass{article}

\usepackage{microtype}
\usepackage{graphicx}
\usepackage{subcaption}
\usepackage{booktabs}
\usepackage{tabularx}
\usepackage{hyperref}
\usepackage{amsmath}
\usepackage{amssymb}
\usepackage{multirow}
\usepackage{array}
\usepackage{xcolor}
\usepackage{balance}

\newcommand{\theHalgorithm}{\arabic{algorithm}}
\newcommand{\norm}[1]{\left\lVert #1 \right\rVert}
\newcommand{\R}{\mathbb{R}}
\newcommand{\od}{\mathbin{\odot}}
\newcommand{\DAWN}{DAWN-SRW}
\newcommand{\op}{\mathcal{O}}
\newcommand{\smallhead}[1]{\vspace{0.25em}\noindent\textbf{#1}}
\newcolumntype{Y}{>{\raggedright\arraybackslash}X}

\usepackage[preprint]{icml2026}
\makeatletter
\renewcommand{\Notice@String}{}
\makeatother
\hypersetup{pdfsubject={}}

\icmltitlerunning{\smash{DAWN-SRW: Learned State Transitions}}

\begin{document}

\twocolumn[
\icmltitle{DAWN-SRW: A Computational Model of Learned State Transitions}

\begin{icmlauthorlist}
\icmlauthor{Seungho Choi}{ind}
\end{icmlauthorlist}

\icmlaffiliation{ind}{Independent Researcher}
\icmlcorrespondingauthor{Seungho Choi}{madst0614@gmail.com}
\icmlkeywords{Language Modeling, Learned State Transitions, Read-Write Operators, Conditional Computation, Mechanistic Analysis}

\vskip 0.3in
]

\printAffiliationsAndNotice{}

\begin{abstract}
Dense language models operate over learned high-dimensional states, but the
operations producing their transitions and the conditions governing their use
remain largely implicit. This motivates a broader computational question:
\emph{Can computation over learned high-dimensional states be organized as an
explicit, learnable system of state-conditioned control over operations?}
\DAWN{} tests this question in language modeling through a narrower architectural
hypothesis: transition-producing language-model computation can be factorized
into identifiable learned operations and explicit state-conditioned control that
determines how those operations contribute to the subsequent state transition.
Here, identifiable means
referable as a stable architectural object and available for measurement or
intervention, not necessarily human-interpretable. The evaluated v4172 model is
one concrete realization. It
uses persistent rank-1 read-write (RW) banks, learned operation-address control,
continuous-valued applicability coefficients with exact support, and normalized
additive composition to construct Q/K/V attention interfaces and direct RST
updates; operator visibility remains global and RW execution remains dense. At
393.8M parameters and 39.999B tokens, \DAWN{} retains performance against a
matched dense Transformer: C4 cross-entropies are 2.9402 versus 2.9359, and
six-task zero-shot means are 0.4629 versus 0.4597. A descriptive 1.272B
evaluation reaches 2.8359 cross-entropy and a 0.4858 zero-shot mean. The model
exhibits exact input-conditioned applicability support, and a discovery-frozen
4,096-site Q/K-centered operator circuit passes suppression, four
matched-control families, and exact restoration on independent validation and
held-out indirect-object-identification inputs. Fail-closed results delimit where
localization did not succeed. Together, these results provide a constructive
test of the narrower architectural hypothesis under the evaluated conditions
and evidence relevant to the broader computational question, without
establishing human-readable semantics, general editability, benchmark
superiority, or hardware efficiency.
\end{abstract}

\section{Introduction}
\label{sec:introduction}

Modern language models achieve strong performance by repeatedly transforming
learned high-dimensional representations, but those transformations do not
expose the constituent operations or the conditions under which a state uses
them. The long-term motivation is to make learned computation available as
explicit computational objects that can be identified, analyzed, and causally
intervened on. Such an interface could eventually support semantic
interpretation, reliable editing, or more programmable learned systems. Those
outcomes are motivations, not claims established by this paper.

We begin from a broader computational question:
\begin{quote}
\emph{Can computation over learned high-dimensional states be organized as an
explicit, learnable system of state-conditioned control over operations?}
\end{quote}
At this level, a learned representation \(x\) is the current computational state,
the state determines an applicability profile
\(C(x)=(c_1(x),\ldots,c_N(x))\) that represents state-conditioned control, and
controlled operation contributions construct the computation that produces the
subsequent state transition:
\[
x\longrightarrow C(x)\longrightarrow
\text{controlled operations}\longrightarrow x'.
\]
Control answers, at the architectural level, which learned operations are
applicable to the current state and to what degree; the operations determine
what computation is performed on that state. \DAWN{} studies this broad question
in the concrete domain of language modeling.

At a high level, conventional programming makes computation explicit through
states, operations, and rules governing their execution. Operational accounts
describe execution through transitions between configurations
\cite{plotkin1981structural}:
\[
\begin{gathered}
\text{current state}
\longrightarrow
\text{conditions / applicable operations}\\[-0.15em]
\longrightarrow
\text{execution}
\longrightarrow
\text{next state}.
\end{gathered}
\]
Such explicit structure supplies objects through which computation can be
referred to and analyzed. We use this program-like organization only in the
limited sense shown above---state, conditions and applicable operations,
execution, and next state. It does not imply symbolic programs, discrete or
human-readable instructions or predicates, or a formal operational semantics
for \DAWN{}.

Language models likewise evolve computational states by repeatedly transforming
representations across depth. Let \(X_\ell\) denote the sequence representation
at layer \(\ell\); at a coarse level, a block implements
\[
X_\ell\longrightarrow X_{\ell+1}.
\]
We likewise use this state-evolution view only as an organizing observation,
not as a contribution or a claim about an intrinsic language-model semantics.

In conventional Transformers, the mechanisms producing these state changes
remain largely implicit in dense parameterized transformations. Attention and
feed-forward layers expose module boundaries, but they do not natively expose a
referable set of operations together with explicit state-conditioned control
over their application.

We therefore study a narrower architectural hypothesis:
\begin{quote}
\emph{Can transition-producing language-model computation be factorized into
identifiable learned operations and explicit state-conditioned control that
determines how those operations contribute to the subsequent state transition?}
\end{quote}
This hypothesis separates what computation is available from when and how it is
applied. ``Identifiable'' means that an operation is referable as a stable
architectural object whose contribution can be measured or intervened on; it
does not specify how broadly or how long that identity must persist, and does
not imply a human-readable concept, a discrete instruction, or a uniquely
correct decomposition. The paper provides a constructive architectural test,
not evidence that this factorization is intrinsic to all language models.

The evaluated \DAWN{} v4172 architecture instantiates one point in this broader
design space. It chooses a rank-1 read-write transformation as its RW primitive
and normalized additive composition as its application semantics. For a token
state \(x\),
\begin{align}
\op_i(x)&=\langle x,r_i\rangle w_i,
\label{eq:intro_operator}\\
F(x)&=\frac{1}{Z(x)}\sum_i c_i(x)\op_i(x),
\qquad c_i(x)\in[0,1].
\label{eq:intro_field}
\end{align}
The read direction \(r_i\) determines what is extracted from the current
representation, the write direction \(w_i\) determines the direction of the
primitive's output contribution, and \(c_i(x)\) determines whether and to what
degree that primitive applies. Thus, this realization parameterizes
transition-producing computation through explicit conditional composition of RW
primitives.

The control abstraction is also explicit:
\begin{equation}
x\longmapsto C_a(x)
=\big(c_{a,1}(x),\ldots,c_{a,N_a}(x)\big).
\label{eq:intro_control}
\end{equation}
These components instantiate the narrower hypothesis: the representation
supplies the current state, \(C_a(x)\) controls primitive application, and the
composed result supplies computation used in the next update. Continuous-valued
coefficients with exact support and normalized addition are v4172 design
choices, not requirements of the architectural hypothesis.

More specifically, v4172 uses persistent type-specific QK, V, and RST RW banks
across depth; a state-derived query and threshold \(\tau\); read/write-derived
primitive addresses in a learned spherical geometry; exact applicability
support with continuous-valued coefficients; normalized additive composition;
and deployment across Q, K, V, and the direct residual-state transition (RST).
Each successive state recomputes applicability over the corresponding bank via
\[
x\longmapsto\big(q_a(x),\tau_a(x)\big)
\longmapsto c_{a,i}(x),
\]
using primitive addresses \(k_{a,i}\). The operation-address space implements
control rather than RW execution. The current model evaluates applicability with
global operator visibility and applies the RW banks densely, even when exact
applicability support is a strict subset of the pool.

This realization builds on the read-write view of dense feed-forward
computation, where first-layer directions detect input patterns and second-layer
directions contribute outputs \cite{geva2021transformer}. That analysis reveals
functional directions inside an existing dense layer. \DAWN{} instead turns RW
factorization into an architectural constraint: reusable RW units are
executable RW primitives whose state-dependent contributions are
directly controlled and observable.

The same formulation produces the query, key, and value states used by causal
self-attention and the direct update applied to the post-attention residual
state.
Q and K produce separate outputs from one shared QK pool; V and the
post-attention residual transition (RST) use separate pools. Operator pools and
applicability projections are reused across depth, while normalization and
attention output projections remain layer-specific.

The empirical study follows a fixed claim hierarchy. First, viability is tested
by comparing a 393.8M parameter \DAWN{} checkpoint with a parameter- and
token-matched dense Transformer on C4 and six zero-shot tasks, with a descriptive
1.272B evaluation. Second, conditional organization is tested by measuring exact
input-conditioned applicability support. Third, causal accessibility is tested
by localizing a discovery-only indirect object identification (IOI) circuit. The circuit is
frozen before confirmatory data are opened, then tested with suppression,
matched controls, and exact restoration on independent validation and fully
held-out test inputs. Fourth, three fail-closed boundaries are reported: RAVEL
localization, ARC causal validation, and 1.3B IOI localization. Efficiency is
not a headline result: the current implementation scores and applies full
operator pools densely.

Our contributions are:
\begin{enumerate}
    \item \textbf{An explicit architectural factorization and one realization.}
    We formulate transition-producing computation in terms of identifiable
    learned operations and state-conditioned control that determines how those
    operations contribute to the subsequent state transition. We instantiate
    that factorization in
    v4172 through persistent rank-1 RW banks, learned spherical operation-address
    control, exact applicability support with continuous-valued coefficients,
    and normalized composition across Q/K/V/RST.
    \item \textbf{A constructive empirical test with bounded evidence.} We test
    viability at 400M scale, report a descriptive 1.3B evaluation, characterize
    conditional organization through applicability support, and test causal
    accessibility through a discovery-frozen IOI circuit confirmed by
    suppression, matched controls, and exact restoration. Fail-closed results
    bound where the same localization program did not succeed.
\end{enumerate}

\section{Related Work}
\label{sec:related}

\subsection{State-dependent and dynamically parameterized transformations}

State-dependent parameterizations have long precedents in nonlinear control
and neural sequence models. State-dependent Riccati equation techniques use
state-dependent coefficient representations in nonlinear control design
\cite{cloutier1997state}; we cite this as conceptual lineage, not as a
mathematical equivalence to \DAWN{}. Multiplicative RNNs allow the current input
to determine the hidden-state transition matrix
\cite{sutskever2011generating}, input-switched affine networks use
input-dependent recurrent weights \cite{foerster2017input}, and HyperNetworks
generate weights for another network \cite{ha2017hypernetworks}. These works
address different settings---feedback control, recurrent dynamics, and generated
neural parameters---but collectively establish that a transformation may depend
on the current state or input. State evolution, state dependence, and
dynamically generated transformations are therefore not contributions of
\DAWN{}. The narrower hypothesis studied here is that transition-producing
language-model computation can be factorized into identifiable learned
operations and explicit state-conditioned control that determines how those
operations contribute to the subsequent state transition.
Persistent RW primitives, address-based control, and normalized additive
composition are choices of the evaluated realization rather than the hypothesis
itself.

\subsection{Conditional computation and fine-grained experts}

Sparsely-gated mixtures of experts activate a routed subset of expert networks
\cite{shazeer2017outrageously}, and Switch Transformers scale this pattern with
simple routing \cite{fedus2022switch}. PEER applies product-key retrieval to
more than a million tiny experts \cite{he2024mixture}; OmniMoE extends this
fine-grained line with vector-level Atomic Experts, a Cartesian Product Router,
and execution scheduling \cite{shi2026omnimoe}. Thus state-dependent routing,
fine-grained experts, large expert pools, and sparse retrieval all precede
\DAWN{}. These systems primarily develop routing, capacity, and efficient sparse
execution for expertized layers. The paper instead studies explicit
state-conditioned operator organization as the base parameterization of
transition-producing computation, and exposes the resulting operator
contributions for causal intervention. Its current implementation still
performs a dense address scan and dense, chunked RW application rather than
indexed sparse execution.

\subsection{Shared computation across depth}

Parameter sharing across depth and layer-agnostic expert reuse also have direct
precedents. Universal Transformers recur over shared layer parameters
\cite{dehghani2019universal}, while MoEUT develops a shared-layer Transformer
with expertized feed-forward and attention computation
\cite{csordas2024moeut}. More recent MoE designs move expert capacity beyond
layer-local ownership: MoUE reuses a universal layer-agnostic expert pool across
layers \cite{chen2026mixture}, and UniPool exposes one global pool through
independent per-layer routers \cite{huang2026unipool}. The former line studies
recurrence or shared-layer computation, while the latter relaxes layer-local
expert ownership. \DAWN{} therefore does not claim depth sharing or a global
pool as new. It persistently reuses fine-grained, type-specific QK, V, and RST
RW banks; each depth's current representation recomputes applicability over the
same bank, and the shared operation capacity constructs Q/K/V attention
interfaces or a direct residual transition.

\subsection{Conditional attention and unified expertized transformations}

Conditional computation has also been applied beyond Transformer feed-forward
layers. SwitchHead routes among attention-side experts and reduces the number
of attention matrices evaluated \cite{csordas2024switchhead}. UMoE reformulates
attention to expose an FFN-like expert structure, then studies a shared expert
set across attention and FFN computation \cite{yang2025umoe}. They establish
that expertized transformations and parameter sharing need not be confined to
the FFN. These works rule out positioning attention-side experts or conditional
attention itself as a \DAWN{} contribution. \DAWN{} instead applies one
state-conditioned RW operator-composition abstraction consistently to the
learned Q, K, and V interfaces and the direct post-attention residual-state
transition, with persistent type-specific RW banks reused across depth.

\subsection{Read-write primitives and memory views}

Transformer feed-forward layers admit a key-value-memory interpretation in
which read-side directions detect input patterns and write-side directions
contribute outputs \cite{geva2021transformer}. Product Key Memory factorizes
retrieval keys \cite{lample2019large}, and scalable memory layers extend large
content-addressable trainable memories \cite{berges2024memory}. Fast Weight
Programmers provide an algebraic and programming precedent based on additive
outer-product updates \cite{schlag2021linear}. MoRAM explicitly treats rank-1
adapters as associative-memory experts and uses their intrinsic keys for
self-activation \cite{lu2026little}. Together these lines establish separate
precedents for read/write interpretation, scalable retrieval, outer-product
algebra, and rank-1 memory experts. RW, rank-1, and memory views are therefore
not new individually. \DAWN{} uses RW factorization as an executable
architectural primitive rather than only a post-hoc interpretation: the
architecture directly exposes each primitive's state-dependent applicability,
executed response, and contribution to the composed computation.

\subsection{Neural programming and programming interpretations}

Neural Programmer-Interpreters represent and execute programs from supervised
traces \cite{reed2016npi}, while Neural Programmer differentiably selects
operations and data \cite{neelakantan2016neuralprogrammer}. Fast Weight
Programmers connect outer-product computation to elementary instructions
\cite{schlag2021linear}. \DAWN{} does not induce a symbolic program, define a
formal operational semantics, or require individual operators to be
semantically interpretable. Here ``program-like'' refers only to explicit
learned operations, state-conditioned applicability, and their controlled
application in producing state transitions. Plotkin's operational perspective
is an organizational reference rather than a formal semantics for the model.

\subsection{Mechanistic analysis and causal intervention}

Mechanistic interpretability studies behavior through localized components and
causal interventions. IOI work isolates an attention-head circuit and evaluates
faithfulness, completeness, and minimality
\cite{wang2023interpretability}. AtP* provides a scalable approximation to
activation patching for localizing behavior to model components
\cite{kramar2024atp}. RAVEL shows why neuron-level representations may remain
entangled and evaluates methods against multiple causal criteria
\cite{huang2024ravel}. This literature distinguishes a candidate localization
from evidence that an intervention causally controls behavior. \DAWN{} supplies
precise RW-operator intervention targets, but explicit units do not guarantee
semantic interpretability and localization alone is not causal confirmation.
The present evidence therefore preserves a discovery freeze followed by
suppression, matched controls, and exact same-example restoration, and does not
generalize the IOI result beyond the specified task, model, and intervention.

Across these lines, state dependence, conditional routing, depth sharing,
attention-side experts, rank-1 RW structure, and programming interpretations all
have clear precedents. The contribution here is a constructive architectural
test of whether such an organization can be explicitly instantiated in language
modeling. The narrower tested factorization is instantiated consistently across
Q/K/V/RST and evaluated for viability, conditional organization, and causal
accessibility.

\section{DAWN-SRW}
\label{sec:method}

\subsection{State Evolution as the Starting Abstraction}
\label{sec:state-evolution}

Let
\[
X_\ell=(x_{\ell,1},\ldots,x_{\ell,T})\in\R^{T\times D}
\]
denote the sequence representation at depth \(\ell\). At a coarse level, a
Transformer-like block evolves this state through a relation-dependent
attention update followed by a transformation of the contextualized
representation:
\begin{align}
Y_\ell&=X_\ell+A_\ell(X_\ell),\\
X_{\ell+1}&=Y_\ell+R_\ell(Y_\ell).
\label{eq:state_evolution}
\end{align}
This decomposition only fixes the computational abstraction and notation.
\DAWN{} concerns how the transformations producing these updates are
parameterized.

\subsection{Explicit State-Conditioned Operator Parameterization}
\label{sec:transition-hypothesis}

The architectural hypothesis is broader than any particular primitive, control
geometry, or application rule. It separates identifiable learned operations
from explicit state-conditioned control that determines how those operations
contribute to the subsequent state transition. The evaluated v4172
architecture instantiates that hypothesis through reusable RW primitives, an
applicability profile, and normalized additive composition. For a token state
\(x\in\mathcal{S}=\R^D\), let
\(a\in\{Q,K,V,\mathrm{RST}\}\) index the computation type and define
\begin{equation}
F_a:\mathcal{S}\rightarrow\R^D.
\label{eq:composed_output}
\end{equation}
The Q, K, and V outputs form the state-conditioned interfaces through which
relational interaction is computed, whereas \(F_{\mathrm{RST}}\) directly
specifies a residual-state update. Each composed output is explicitly
decomposed into contributions from reusable RW primitives under state-dependent
control. This RW factorization, its rank-1 granularity, the operation-address
coupling, spherical geometry, and normalized additive composition are design
choices of v4172. They provide one concrete test of the hypothesis rather than
defining it.

At the computational level, this view separates two questions for a state
\(x\): which learned operations are applicable, and what transformations those
operations perform. \DAWN{} makes both sides architecturally explicit:
\(C_a(x)\) represents state-conditioned control over primitive application,
whereas the RW primitives represent the transformations whose contributions are
controlled. Spherical operation-address geometry realizes this applicability
profile in v4172; it does not define the abstract control representation or the
broader hypothesis.

\begin{figure*}[t]
\centering
\setlength{\fboxsep}{4pt}
\small
\textbf{Language modeling as representation-state evolution}\\[0.35em]
\begin{tabular}{c@{\hspace{0.12em}\(\longrightarrow\)\hspace{0.12em}}c@{\hspace{0.12em}\(\longrightarrow\)\hspace{0.12em}}c@{\hspace{0.12em}\(\longrightarrow\)\hspace{0.12em}}c@{\hspace{0.12em}\(\longrightarrow\)\hspace{0.12em}}c}
\fbox{\parbox[c][0.92cm][c]{0.10\textwidth}{\centering
\textbf{Sequence state}\\\(X_\ell\)}} &
\fbox{\parbox[c][0.92cm][c]{0.17\textwidth}{\centering
\textbf{Relational interaction}\\attention via \(F_Q,F_K,F_V\)}} &
\fbox{\parbox[c][0.92cm][c]{0.10\textwidth}{\centering
\textbf{Contextual state}\\\(Y_\ell\)}} &
\fbox{\parbox[c][0.92cm][c]{0.16\textwidth}{\centering
\textbf{Residual-state transition}\\via \(F_{\mathrm{RST}}\)}} &
\fbox{\parbox[c][0.92cm][c]{0.10\textwidth}{\centering
\textbf{Next state}\\\(X_{\ell+1}\)}}
\end{tabular}\\[-0.05em]
\begin{tabular}{cc}
\parbox{0.34\textwidth}{\centering
\(F_Q,F_K,F_V\) parameterize relational interaction\\\(\Downarrow\)} &
\parbox{0.34\textwidth}{\centering
\(F_{\mathrm{RST}}\) directly updates the residual state\\\(\Downarrow\)}
\end{tabular}\\[0.15em]
\textbf{\DAWN{}: state-conditioned operator composition}\\[0.35em]
\begin{tabular}{c@{\hspace{0.5em}}c@{\hspace{0.5em}}c@{\hspace{0.5em}}c@{\hspace{0.5em}}c}
&& \fbox{\parbox[c][0.68cm][c]{0.23\textwidth}{\centering
\textbf{Applicability control}\\\(C_a(x)=(c_{a,1}(x),\ldots,c_{a,N_a}(x))\)}} && \\[-0.28em]
\fbox{\parbox[c][0.68cm][c]{0.10\textwidth}{\centering
\textbf{Token state}\\\(x\)}} & \(\nearrow\) && \(\searrow\) &
\fbox{\parbox[c][0.68cm][c]{0.18\textwidth}{\centering
\textbf{Normalized composition}\\\(F_a(x)\)}} \\[-0.28em]
& \(\searrow\) && \(\nearrow\) & \\[-0.28em]
&& \fbox{\parbox[c][0.86cm][c]{0.23\textwidth}{\centering
\textbf{Persistent RW bank}\\\(\{\op_{a,i}(x)\}_{i=1}^{N_a}\)\\[-0.12em]
\scriptsize type-specific; shared across depth}} &&
\end{tabular}\\[0.15em]
\caption{The two-level construction. Top: a language-model block evolves a
sequence representation through relational interaction and a subsequent
residual-state transition. Bottom: \(x\) separately determines the applicability
profile and input-dependent responses of a persistent, type-specific RW bank
shared across depth; the two meet in normalized composition \(F_a(x)\). In
v4172, \(q_a(x),\tau_a(x),k_{a,i}\)
implement applicability in address space; RW execution remains in the full
\(D\)-dimensional representation space.}
\label{fig:architecture}
\end{figure*}

\subsection{Read-Write Computational Primitives}
\label{sec:rw}

Each primitive owns learned read and write parameters. With
\(\operatorname{unit}(z)=z/\max(\norm{z}_2,\epsilon)\), define normalized
directions \(\bar r_{a,i}=\operatorname{unit}(r_{a,i})\) and
\(\bar w_{a,i}=\operatorname{unit}(w_{a,i})\). The primitive executes
\begin{equation}
\op_{a,i}(x)=\langle x,\bar r_{a,i}\rangle\bar w_{a,i}
=\bar w_{a,i}\bar r_{a,i}^{\top}x.
\label{eq:rw_operator}
\end{equation}
This is an explicit rank-1 RW primitive. It reads a scalar from the current
state and produces a directed output contribution. The read direction
specifies the extracted scalar, while the write direction specifies the
direction of that contribution. For Q, K, and V, these contributions compose
the interfaces used by attention; for RST, they compose a direct residual-state
update. Unlike a static retrieved value, an RW primitive remains an
input-dependent transformation, with its output varying with the current state
through the read response. The state-to-state transition
is completed only when the composed output participates in the
attention-mediated or residual update.

\subsection{State-Conditioned Applicability}
\label{sec:applicability}

At the abstract level, control is represented by an applicability profile over
reusable RW primitives. For each computation type
\(a\in\{Q,K,V,\mathrm{RST}\}\), the current state induces a profile over all
\(N_a\) primitives,
\begin{equation}
C_a(x)=\big(c_{a,1}(x),\ldots,c_{a,N_a}(x)\big),
\label{eq:applicability_profile}
\end{equation}
whose coordinates specify which primitives participate and to what degree,
rather than assigning the state to a single discrete branch. The evaluated
v4172 realization has exact zero applicability outside learned boundaries and
continuously positive applicability inside them. This is not an algebraic
decomposition of the representation vector; its coordinates control the
contributions of the corresponding RW primitives composed below.

In v4172, normalized coordinates \(q_a(x)\) and \(k_{a,i}\) compare the current
state with primitive \(i\)'s learned address;
Section~\ref{sec:operation-space} gives their construction. A separate direct
projection produces the state-dependent DirectTau boundary:
\begin{align}
\tau_a(x)&=-1+2\sigma(xu_a+\beta_a),\\
\rho_{a,i}(x)&=q_a(x)^\top k_{a,i},\\
m_{a,i}(x)&=\rho_{a,i}(x)-\tau_a(x).
\label{eq:directtau}
\end{align}
Here \(m_{a,i}(x)\) is the applicability margin. The evaluated
\texttt{linear\_angular} rule defines primitive applicability as
\begin{equation}
c_{a,i}(x)=
\operatorname{clip}\!\left(
\frac{m_{a,i}(x)}
{\max(1-\tau_a(x),10^{-4})},0,1\right).
\label{eq:angular_gate}
\end{equation}
We call \(c_{a,i}(x)\in[0,1]\) the \emph{continuous-valued applicability
coefficient} for primitive \(i\) at state \(x\). Its exact state-dependent
applicability support is
\begin{equation}
\mathcal{I}_a(x)=\{i:\rho_{a,i}(x)>\tau_a(x)\}
=\{i:c_{a,i}(x)>0\}.
\label{eq:applicable_support}
\end{equation}
Outside this set, a primitive contributes exactly zero; inside it,
\(0<c_{a,i}(x)\leq1\). A primitive's contribution therefore enters
continuously from zero at the learned boundary. This combines state-dependent
sparse applicability support with a continuous-valued applicability coefficient rather than
merely assigning a dense soft mixture over every primitive.

\subsection{State-Conditioned Operator Composition}
\label{sec:operator-composition}

Applicable primitives compose into the output \(F_a(x)\):
\begin{align}
Z_a(x)&=\left[\max\!\left(\sum_i c_{a,i}(x),1\right)\right]^{p_a},\\
s_a&=\sqrt{D/L},\nonumber\\
F_a(x)&=\frac{s_a}{Z_a(x)}
\sum_i c_{a,i}(x)\op_{a,i}(x).
\label{eq:implemented_composition}
\end{align}
Equivalently,
\begin{equation}
\begin{aligned}
F_a(x)&=A_a(x)x,\\
A_a(x)&=\frac{s_a}{Z_a(x)}
\sum_i c_{a,i}(x)\bar w_{a,i}\bar r_{a,i}^{\top}.
\end{aligned}
\label{eq:state_dependent_operator}
\end{equation}
This form makes the architectural hypothesis explicit. The model does not apply
one fixed dense transformation to every state. Instead, the current
representation determines continuous-valued applicability coefficients, which
construct a state-dependent transformation \(A_a(x)\) from reusable rank-1 RW
primitives. The resulting computation is both input-conditioned and explicitly
decomposed into identifiable operator contributions.

``Continuous-valued'' refers to the applicability coefficient, not to
continuous-time dynamics or a continuous-depth model; layer evolution remains
discrete. The scalar mapping from normalized applicability margin to
applicability is continuous and piecewise linear, with nondifferentiable support
boundaries.

The denominator powers are \(p_Q=p_K=0.5\), \(p_V=1.0\), and
\(p_{\mathrm{RST}}=1.2\). The fixed composition scale is
\(s_a=\sqrt{D/L}\), with no
learned pool-strength parameter.

\subsection{Operation-Address Space}
\label{sec:operation-space}

The operation-address space \(\mathcal{A}=\R^R\) is where primitive
applicability is evaluated, not where RW computation itself is executed. This
addressing construction is one evaluated realization of the state-conditioned
applicability abstraction above. The present experiments do not establish that
this particular realization is necessary or preferable to alternative
state-conditioned control mechanisms. In v4172, \(R=256\), and a normalized
query is a direct projection of the current input state:
\begin{equation}
q_a(x)=\operatorname{unit}(xP_a+b_a)\in\mathcal{A}.
\label{eq:query}
\end{equation}
A shared probe pair \(P_r,P_w\in\R^{D\times R}\) maps all operator pools'
functional read and write directions into address features. The
operator address is
\begin{align}
u_{a,i}&=\operatorname{unit}(\bar r_{a,i}^{\top}P_r),&
v_{a,i}&=\operatorname{unit}(\bar w_{a,i}^{\top}P_w),\nonumber\\
k_{a,i}&=\operatorname{unit}(u_{a,i}\od v_{a,i})\in\mathcal{A}.
\label{eq:rw_address}
\end{align}
This generalized coordinate-wise bilinear mapping makes address identity a
learned function of what a primitive reads and where it writes, rather than an
independent expert ID. It is the addressing realization used to compute the
applicability coefficients in Equation~\ref{eq:angular_gate}, not the paper's
definition of a state transition. Shared probes align QK, V, and RST addresses
to one coordinate convention, while their operator rows remain in QK-, V-, and
RST-specific pools.

Because both \(q_a(x)\) and \(k_{a,i}\) are unit-normalized, they lie on the
address sphere \(\mathbb{S}^{R-1}\). Let
\(\theta_{a,i}(x)=\arccos\!\big(q_a(x)^\top k_{a,i}\big)\) and
\(\theta_a(x)=\arccos\tau_a(x)\). The applicability support is therefore
\[
\mathcal{I}_a(x)
=\{i:\theta_{a,i}(x)<\theta_a(x)\},
\]
a spherical cap centered at the query \(q_a(x)\) with state-dependent angular
radius \(\theta_a(x)\). Geometrically, \(q_a(x)\) specifies the center of the
applicability query, \(\tau_a(x)\) specifies its angular support boundary, and
\(c_{a,i}(x)\) specifies the continuous-valued coefficient with which primitive \(i\)
applies. When the
numerical guard in Equation~\ref{eq:angular_gate} is inactive, an applicable
primitive has
\[
c_{a,i}(x)=
\frac{\cos\theta_{a,i}(x)-\tau_a(x)}{1-\tau_a(x)},
\]
which rises continuously from zero at the cap boundary to one at its center.
Equivalently, \(c_{a,i}(x)\) is a normalized continuous depth within the cap.
Equation~\ref{eq:angular_gate} is the exact implementation, retaining the
\(10^{-4}\) denominator guard and clipping.

\subsection{Attention-Mediated and Residual Transitions}
\label{sec:block}

For a sequence state \(X_\ell\), layer-specific normalization gives
\(\widetilde X_\ell=\operatorname{LN}^{(1)}_\ell(X_\ell)\). Three
state-conditioned operator compositions produce the attention interfaces:
\begin{align}
Q_\ell&=F_Q(\widetilde X_\ell),\nonumber\\
K_\ell&=F_K(\widetilde X_\ell),\nonumber\\
V_\ell&=F_V(\widetilde X_\ell),\nonumber\\
A_\ell&=\operatorname{softmax}\!\left(
\frac{Q_\ell K_\ell^\top}{\sqrt{d_h}}+M_{\mathrm{causal}}\right)
\!,\nonumber\\
\Delta X_{\mathrm{attn},\ell}&=A_\ell V_\ell W^O_\ell,\\
Y_\ell&=X_\ell+\Delta X_{\mathrm{attn},\ell}.
\label{eq:attention}
\end{align}
Q and K produce separate outputs from the same 4,600-row QK pool; V produces an
output from a separate 14,500-row pool. After attention, a second layer-specific
normalization \(\widetilde Y_\ell=\operatorname{LN}^{(2)}_\ell(Y_\ell)\) is
evaluated anew using the 29,300-row RST operator pool:
\begin{equation}
X_{\ell+1}=Y_\ell+F_{\mathrm{RST}}(\widetilde Y_\ell).
\label{eq:rst}
\end{equation}
The block thus applies an attention-mediated state transition followed by a
direct residual-state transition. RST applicability depends on the
contextualized post-attention state, not the pre-attention query.

\subsection{Cross-Depth Operator Reuse}
\label{sec:reuse}

The QK, V, and RST read/write pools; the two RW-to-address probes; the Q/K/V
query projection; the RST query projection; and their DirectTau projections
are shared across all \(L\) layers. Q and K share read/write rows but retain
separate query slices, tau values, applicability coefficients, and outputs. The
layer-specific parameters are two LayerNorms and the attention output
projection \(W^O_\ell\). Positional and token embeddings and the final
normalization are model-level parameters. Consequently, \DAWN{} reuses
RW primitives across depth without claiming complete block sharing.

\subsection{Implemented Addressing and Compute Boundary}
\label{sec:implemented-boundary}

For computation type \(a\), define \emph{operator visibility} \(\mathcal{V}_a(x)\) as the
set of operator identities made available to the control mechanism for
applicability evaluation at state \(x\). Visibility denotes the candidate range
available to that mechanism, not whether operators are observable to a human
analyst. In the evaluated v4172 realization, every operator in the type-specific
pool is visible:
\[
\mathcal{I}_a(x)\subseteq\mathcal{V}_a(x)=\{1,\ldots,N_a\},
\]
with \(\mathcal{I}_a(x)\) generally a strict subset.

The implementation shards operator rows across the model mesh and processes
them in dense chunks. The 400M configuration uses one QK chunk, one V chunk,
and four RST chunks. Every query scores the full local address pool; every
operator row participates in dense read and write contractions before
cross-shard reduction. Exact applicability support can be sparse while operator
visibility remains global and RW execution remains dense; it is not an indexed
sparse kernel.

The evaluated architecture keeps RW execution in the full residual dimension
\(D\); the \(R\)-dimensional operation-address space is used only for
applicability evaluation. An executor could omit exactly zero RW terms without
changing the mathematical output, but such support-restricted RW execution is
not implemented or timed here.

\section{Experimental Setup}
\label{sec:setup}

\subsection{Models and training}

Table~\ref{tab:model-spec} summarizes the evaluated checkpoints. The primary
\DAWN{} model has 393,800,708 measured parameters, operation-address dimension
\(R=256\), and 48,400 operator rows across QK, V, and RST. It was trained for
39,999,504,384 packed C4 tokens with sequence length 512 and global training
batch 1,024. Both dropout rates are zero.

\begin{table*}[t]
\centering
\caption{Model and training specification. The 1.3B row is descriptive and
does not form a controlled scaling curve with the 400M rows.}
\label{tab:model-spec}
\small
\resizebox{0.99\textwidth}{!}{
\begin{tabular}{lrrrrrrrr}
\toprule
Model & Parameters & Tokens & \(D\) & Layers & Heads & QK / V / RST rows & Checkpoint\\
\midrule
\DAWN{} 400M & 393,800,708 & 39,999,504,384 & 2,304 & 18 & 36 &
4,600 / 14,500 / 29,300 & 76,293\\
Dense Transformer 400M & 393,896,960 & 39,999,504,384 & 1,280 & 18 & 20 &
-- & 76,293\\
\DAWN{} 1.3B & 1,272,038,404 & 19,999,981,568 & 4,096 & 24 & 32 &
8,534 / 26,904 / 54,352 & 87,193\\
\bottomrule
\end{tabular}
}
\end{table*}

\subsection{Matched 400M dense control}

The dense control has 393,896,960 measured parameters, width 1,280, feed-forward
width 5,120, 18 layers, and 20 heads. It uses the same tokenizer, vocabulary,
sequence length, training-token count, seed, batch size, mesh, learning rate,
warmup ratio, weight decay, and zero dropout as the primary model. The
parameter difference is 96,252, or 0.0244\% of the dense model. This comparison
tests performance retention under matched scale and training exposure; it is
not a formal equivalence test.

\subsection{Evaluation protocols}

The C4 protocol evaluates 610 complete global batches from the fixed
10M-token capped packed-validation view: 19,520 sequences and 9,974,720 valid
next-token targets. Cross-entropy, perplexity, and next-token accuracy use
float32 scoring. The zero-shot protocol uses \texttt{lm-eval} 0.4.2 with no
few-shot examples and no task limit on LAMBADA, HellaSwag, PIQA, ARC-Easy,
ARC-Challenge, and WinoGrande. The primary metric is raw accuracy for LAMBADA
and WinoGrande and normalized accuracy for the other four tasks. The mean is
an unweighted mean of these six metrics.

\subsection{Mechanistic localization and frozen intervention protocol}
\label{sec:protocol}

We define a site by layer, computation type, token position, and global operator
identity.
IOI discovery uses 119 independent paired-correct units. Sites are ranked by
absolute discovery mean contribution importance. We choose the smallest
audited prefix that reaches at least 0.70 cumulative absolute importance and
0.95 split top-\(k\) overlap. The resulting 4,096-site Q/K/V/RST/layer
composition and intervention specification are frozen before validation.

The primary causal score is the summed-log-probability margin between the
correct indirect object and the source name. Circuit-wide suppression zeros
only the selected operators' execution numerator at target positions while
retaining the full production denominator. Four control families each contain
100 frozen-size circuits that exclude circuit sites: equal-count random,
same-layer random, activation-matched, and route-frequency-matched across
Q/K/V/RST. Exact
restoration reinserts the same-example selected numerator after suppression,
again under the production denominator. Validation uses 123 independent
paired-correct units; the held-out test accessor is opened only after the
validation result and unchanged specification are finalized. We report 95\%
bootstrap confidence intervals and two-sided paired permutation tests where
specified.

\section{Results}
\label{sec:results}

\subsection{Language-model performance retention at 400M}

Table~\ref{tab:c4} reports the locked C4 evaluation. \DAWN{} differs from the
dense control by \(+0.004318918\) cross-entropy, \(+0.081537182\) perplexity,
and \(-0.002188432\) accuracy (\(-0.2188\) percentage point). These small
measured gaps support performance retention, not superiority or formal
equivalence.

\begin{table}[t]
\centering
\caption{Packed C4 validation on the identical 9,974,720 target tokens.}
\label{tab:c4}
\small
\begin{tabular}{lrrr}
\toprule
Model & CE & PPL & Accuracy\\
\midrule
\DAWN{} 400M & 2.940213 & 18.919873 & 0.424912\\
Dense 400M & 2.935894 & 18.838336 & 0.427100\\
\DAWN{} 1.3B & 2.835879 & 17.045382 & 0.437099\\
\bottomrule
\end{tabular}
\end{table}

Table~\ref{tab:zero-shot} gives task-level zero-shot results with standard
errors. The matched 400M means are 0.462872 and 0.459652, a difference of
\(+0.003220495\) (\(+0.3220\) percentage point). Task differences vary in
sign, so the mean does not establish benchmark superiority. LAMBADA
perplexities are \(24.0774\pm0.8005\) for \DAWN{} 400M,
\(24.0042\pm0.8073\) for dense 400M, and \(17.6358\pm0.5390\) for
\DAWN{} 1.3B.

\begin{table*}[t]
\centering
\caption{Zero-shot accuracy. Parentheses contain standard errors. HS, AE, and
AC denote HellaSwag, ARC-Easy, and ARC-Challenge.}
\label{tab:zero-shot}
\scriptsize
\resizebox{\textwidth}{!}{
\begin{tabular}{lrrrrrrr}
\toprule
Model & LAMBADA & HS & PIQA & AE & AC & WinoGrande & Mean\\
\midrule
\DAWN{} 400M &
0.352028 (0.006654) & 0.458375 (0.004972) & 0.712187 (0.010563) &
0.484007 (0.010255) & 0.256826 (0.012767) & 0.513812 (0.014047) & 0.462872\\
Dense 400M &
0.353192 (0.006659) & 0.457080 (0.004971) & 0.698041 (0.010712) &
0.475589 (0.010248) & 0.249147 (0.012639) & 0.524862 (0.014035) & 0.459652\\
\DAWN{} 1.3B &
0.395886 (0.006813) & 0.510157 (0.004989) & 0.723069 (0.010440) &
0.509680 (0.010258) & 0.268771 (0.012955) & 0.507498 (0.014051) & 0.485844\\
\bottomrule
\end{tabular}}
\end{table*}

\subsection{Descriptive larger-scale evaluation at 1.3B}

The 1.272B checkpoint reaches 2.835879 C4 cross-entropy, 17.045382
perplexity, 0.437099 token accuracy, and a 0.485844 six-task mean. It was
trained for 19.999B tokens, about half the 400M training-token count, and no
completed matched dense 1.3B result is admitted to this study. These results
therefore describe a larger trained \DAWN{} model; they are neither a
controlled scaling curve nor a dense-model comparison.

\subsection{Input-conditioned RW applicability support}
\label{sec:support}

At all 9,994,240 packed-C4 positions, Q, K, V, and RST each have strict-subset
applicability support at positive applicability margin (Table~\ref{tab:support}).
Mean exact applicability support
is 487.896 Q operators, 386.722 K operators, 812.388 in the shared-pool Q/K
union, 1,649.994 V operators, and 2,017.703 RST operators. This is exact
input-conditioned applicability support under Equation~\ref{eq:angular_gate}.

\begin{table}[t]
\centering
\caption{Exact applicability support at applicability margin \(>0\). Q/K union
counts an operator in the applicability support for Q or K at the same layer and
token.}
\label{tab:support}
\small
\begin{tabular}{lrrr}
\toprule
Computation & Pool & Mean support & Pool fraction\\
\midrule
Q & 4,600 & 487.896 & 10.6064\%\\
K & 4,600 & 386.722 & 8.4070\%\\
Q/K union & 4,600 & 812.388 & 17.6606\%\\
V & 14,500 & 1,649.994 & 11.3793\%\\
RST & 29,300 & 2,017.703 & 6.8864\%\\
\bottomrule
\end{tabular}
\end{table}

Applicability support does not make the present kernel efficient. Static forward
accounting at batch 1 and sequence length 512 is 4.741260116 TFLOPs for
current dense \DAWN{}, 0.818230250 TFLOPs for dense addressing plus RW
application restricted to exact applicability support, and 0.386547057 TFLOPs
for the matched dense Transformer. The corresponding opportunity is 17.2577\% of current
\DAWN{} accounting but still 2.1168 times the dense Transformer
(Figure~\ref{fig:flops}). It is a mathematical conditional-RW opportunity,
not measured latency, energy, or hardware speedup.

\begin{figure}[t]
\centering
\includegraphics[width=\linewidth]{figures/flops_stacked.pdf}
\caption{Static backbone accounting at batch 1 and sequence length 512.
Address scoring remains dense in the estimate restricted to exact applicability support. Current execution
also applies all full-dimensional RW rows. No production latency claim is made.}
\label{fig:flops}
\end{figure}

\subsection{IOI circuit localization}

Discovery produces split-half rank stability 0.9139788 over 456,697 ranked
sites. The smallest prefix passing both frozen discovery gates contains 4,096
sites and 0.719249 cumulative absolute importance. It spans all 18 layers and
contains 1,744 Q sites, 2,350 K sites, two V sites, and no RST sites. Split
top-\(k\) overlap is 0.974854 (Jaccard 0.950941). These statistics identify a
Q/K-centered candidate circuit; localization alone is not a causal claim.

\subsection{Validation and held-out causal intervention}
\label{sec:causal}

Table~\ref{tab:ioi} shows confirmatory results. On validation, suppression
moves the mean margin from 3.465456 to \(-4.399331\), a drop of 7.864787
(95\% CI [7.490080, 8.241263]), and accuracy falls from 1.000000 to 0.008130.
The frozen effect exceeds every replicate in all four 100-circuit control
families. The most competitive control is activation-matched, with mean drop
5.400490; the frozen-minus-control difference is 2.464297 (95\% CI
[2.076200, 2.832697], paired permutation \(p=0.00049975\)). Exact restoration
returns accuracy to 0.991870 and recovers 0.999661 of the suppression effect
(95\% CI [0.998655, 1.000584]).

\begin{table*}[t]
\centering
\caption{Frozen IOI circuit results. Accuracy arrows show intact to
suppressed. Recovery is the fraction of the suppression effect recovered by
exact same-example numerator restoration.}
\label{tab:ioi}
\small
\resizebox{0.98\textwidth}{!}{
\begin{tabular}{lrrrrrr}
\toprule
Split & Units & Intact margin & Suppressed margin & Drop (95\% CI) &
Accuracy & Recovery / restored acc.\\
\midrule
Validation & 123 & 3.465456 & -4.399331 &
7.864787 [7.490080, 8.241263] &
1.000000 \(\rightarrow\) 0.008130 & 0.999661 / 0.991870\\
Held-out test & 114 & 3.254720 & -4.155613 &
7.410332 [7.016485, 7.818472] &
1.000000 \(\rightarrow\) 0.008772 & 0.998406 / 1.000000\\
\bottomrule
\end{tabular}
}
\end{table*}

The fully held-out test reuses the unchanged site identity, scoring rule,
intervention, controls, seeds, and thresholds. Suppression moves margin from
3.254720 to \(-4.155613\), a drop of 7.410332 (95\% CI
[7.016485, 7.818472]), while accuracy falls to 0.008772. Again the frozen
effect exceeds every replicate in all four control families. The
activation-matched mean drop is 5.472579; the frozen-minus-control difference
is 1.937754 (95\% CI [1.564818, 2.319706], \(p=0.00049975\)). Restoration
recovers 0.998406 of the effect (95\% CI [0.997324, 0.999385]) and restores
accuracy to 1.000000.

Together, suppression, matched-control separation, and exact restoration
support the claim that this discovery-frozen Q/K-centered operator circuit is
causally necessary for the evaluated IOI behavior on independent validation
and fully held-out inputs. ``Necessary'' is intervention-specific: it does not
mean the sites form a unique or human-readable algorithm.

\subsection{Negative and fail-closed results}
\label{sec:negatives}

The same workflow produces informative negative boundaries
(Table~\ref{tab:negative}). RAVEL passes the behavioral gate but fails
variable-localization stability, so no circuit is frozen and no confirmatory
split is opened. ARC discovery freezes a stable 4,096-site Q/K circuit, but
validation suppression has the wrong sign and fails control separation;
restoration passes, and test remains closed. For 1.3B IOI, behavioral
eligibility and rank stability pass, but the largest allowed 4,096-site prefix
reaches only 0.635144 cumulative importance, below the preregistered 0.70
threshold. No circuit is frozen and no validation or test is opened.

\begin{table*}[t]
\centering
\caption{Fail-closed mechanistic results. Later phases were not opened after a
failed preregistered gate.}
\label{tab:negative}
\small
\begin{tabularx}{\textwidth}{l l Y Y}
\toprule
Study & Passed & Failed boundary & Decision\\
\midrule
RAVEL 400M & Behavioral eligibility & Variable stability:
0.628762 / 0.676243 / 0.660622 & No frozen circuit; no validation, transfer, or test\\
ARC 400M & Stable discovery; 4,096 Q/K sites; restoration &
Suppression effect wrong sign; control separation failed &
Reject causal replication; held-out test closed\\
IOI 1.3B & Behavioral eligibility; rank stability 0.907188 &
4,096-site cumulative importance 0.635144 \(<0.70\) &
No frozen circuit; validation and test closed\\
\bottomrule
\end{tabularx}
\end{table*}

\balance
\section{Discussion}
\label{sec:discussion}

\subsection{What the results establish}

The constructive test supports three ordered claims. First, the matched 400M
comparison establishes viability under the measured protocols. Second, exact
input-conditioned applicability support establishes that control is
state-dependent in the evaluated model, even though operator visibility remains
global and RW execution remains dense. Third, the discovery-frozen IOI result
establishes causal accessibility for the specified circuit: suppressing its
operator contributions changes behavior, matched controls do not reproduce that
effect, and exact
same-example restoration recovers it on independent validation and held-out
data. RAVEL, ARC, and 1.3B IOI are fail-closed boundaries showing that successful
localization is not automatic across tasks or scales.

These results do not establish that language reasoning is a discrete program,
that the factorization is intrinsic or optimal, or that identifiable operations
have human-readable semantics. They show that one explicit factorization can
support language modeling, exhibit state-conditioned organization, and provide
causal intervention targets under the evaluated conditions.

\subsection{From a Computational Principle to an Open Design Space}

If learned high-dimensional computation can be organized as state-conditioned
control over operations, what properties should the conditions, operations,
their coupling, and their execution have? The present results provide one
constructive test in language modeling, not a universal answer. Within this open
design space, the narrower architectural hypothesis separates identifiable
learned operations from explicit state-conditioned control that determines how
those operations contribute to
the subsequent state transition. ``Identifiable'' is
an architectural property: an operation is referable as a stable object that
can be measured and intervened on. It does not prescribe the scope or duration
of that identity and is not a claim of semantic purity. Operations may remain
distributed or polysemantic, and a behavior may depend on many operation sites.

v4172 is one explicit realization: its rank-1 granularity, address-based control,
cross-depth persistence, and normalized additive composition are design choices
rather than definitions of the broader computational principle.

\paragraph{What constitutes a learned condition?}
v4172 represents \(C_a(x)\) through an explicit applicability profile derived in
a learned operation-address geometry. Other systems could couple applicability
directly to functional parameters, use intrinsic read-side activation as in
associative-memory-style systems, or use separate control representations. The
present evidence does not establish that one organization is preferable.

\paragraph{What constitutes an operation?}
Rank-1 RW primitives are one granularity. Operations could be finer, coarser,
structured, or hierarchical, and their identities could persist over different
scopes. Identifiability requires only a referable architectural object available
for measurement or intervention; it does not imply semantic purity, and
operations may remain distributed or polysemantic.

\paragraph{How should control and operation be coupled?}
v4172 derives operation addresses from read/write parameters, coupling
functional identity and applicability control without making them identical.
Other designs could separate the two more strongly or couple applicability more
directly to the transformation parameters.

\paragraph{How should applicable operations compose?}
v4172 uses simultaneous normalized additive composition. Other computational
systems could use sequential, ordered, hierarchical, recurrent, or otherwise
structured composition; none of these alternatives is evaluated here.

\paragraph{How much global operator visibility is necessary?}
The current model exposes the full type-specific operator pool to applicability
evaluation. Whether expressive learned control requires global operator
visibility or can preserve its conditional structure when operator visibility is
candidate-restricted remains an empirical question.

\paragraph{How does logical conditional structure map to physical execution?}
Applicability support, operator visibility, and physical execution are distinct.
The execution implications of this separation are developed below.

\subsection{From Logical Conditional Structure to Physical Execution}
\label{sec:future-visibility}

The present results establish state-conditioned applicability support: for a
given representation state, many operators have exactly zero applicability.
They do not establish candidate-restricted operator visibility or sparse physical
execution. In v4172,
\[
\mathcal{I}_a(x)\subseteq\mathcal{V}_a(x)=\{1,\ldots,N_a\},
\]
so applicability is evaluated over the full type-specific pool before
zero-contribution operators are excluded, and RW application remains dense.

The execution problem separates three levels: applicability support, operator
visibility, and physical execution. Applicability support specifies which
operations contribute; operator visibility specifies which operator identities
the control mechanism can consider when determining applicability; physical
execution asks whether the resulting computation can be organized into
accelerator-efficient execution units. \textbf{Sparse applicability support does
not by itself imply hardware-exploitable sparse computation.} Fine-grained
state-conditioned control may conflict with hardware-efficient execution
granularity, memory locality, dispatch, or communication.

The immediate empirical question is how much operator visibility and executed
computation can be reduced, and at what physical execution granularity, while
preserving language-model quality and the learned conditional structure. More
generally, it remains open whether high-dimensional learned control can retain
broad conditional expressivity while acquiring coarse physical locality suitable
for efficient execution.

\subsection{Limitations}

The present results do not establish that transition-producing computation has
been factorized into an efficient conditional computational structure. Although
each state has strict-subset applicability support, applicability is still
evaluated with global operator visibility over large type-specific pools.
Performance retention may therefore depend partly on the representational
capacity provided by a large distributed operator bank and full-pool applicability
evaluation, rather than on a compact or intrinsically efficient decomposition
of language-model computation.

More fundamentally, the present results do not establish that
condition--operation factorization is a necessary, intrinsic, privileged, or
optimal organization of learned high-dimensional computation. They establish
only that one explicit realization is viable under the evaluated language-model
conditions.

Relatedly, the learned spherical address geometry is not shown to reflect an
intrinsic or compressible organization of language operations. It may instead
serve as a high-capacity substrate in which useful operators are distributed
broadly enough for global state-conditioned access to recover the required
transformations. The present study therefore does not establish how much
operator visibility is required to preserve language-model expressivity, or
whether candidate evaluation and executed computation can be restricted
substantially further.

The current implementation evaluates full address and operator pools densely.
Exact applicability support does not imply current sparse execution, and even
the accounting restricted to exact applicability support remains 2.1168 times
the matched dense Transformer's
backbone accounting. No production latency, energy, or hardware advantage is
measured.

The central causal evidence is confined to IOI and does not automatically
generalize to every task. RAVEL, ARC, and the selected 1.3B IOI replication
provide valid negative boundaries. The 400M and 1.3B checkpoints have
different parameter and training-token budgets, so they do not define a
controlled scaling curve. Operator-level causal units need not map one-to-one
to human-readable concepts. Finally, the present work does not establish language
reasoning as universally discrete operator programs.

\section{Conclusion}

This paper begins from a broader computational question: whether computation
over learned high-dimensional states can be organized as an explicit, learnable
system of state-conditioned control over operations. It tests this question in
language modeling through a narrower architectural hypothesis:
transition-producing computation can be factorized into identifiable learned
operations and explicit state-conditioned control that determines how those
operations contribute to the subsequent state transition.

v4172 is one concrete realization, using persistent rank-1 RW banks,
read/write-derived spherical addresses, exact applicability support with
continuous-valued coefficients, and normalized additive composition across
Q/K/V/RST. The empirical results provide a bounded constructive test: measured
400M viability, descriptive 1.3B evidence, input-conditioned applicability
support, and causal accessibility through a discovery-frozen IOI circuit with
suppression, matched controls, and exact restoration. Fail-closed results mark
where the same localization program did not succeed.

These findings do not establish an intrinsic factorization, human-readable
operator semantics, general editability, benchmark superiority, or current
hardware efficiency. They instead motivate an open research program on how
learned conditions are represented, how operations are defined and identified,
how control and operation should be coupled, how applicable operations compose,
how persistence and hierarchy should be structured, and how logical conditional
structure can be mapped to efficient physical execution. Candidate-restricted
operator visibility and execution are one immediate instance of this broader
program.

\appendix

\section{Exact Addressing and Composition Equations}
\label{app:equations}

This appendix restates the production computation as a compact reference. Raw
read and write parameters are normalized in float32 before being cast for
dense contractions:
\[
\bar r_i=\frac{r_i}{\max(\norm{r_i}_2,\epsilon)},\qquad
\bar w_i=\frac{w_i}{\max(\norm{w_i}_2,\epsilon)}.
\]
One pair of model-level probes is shared across QK, V, and RST:
\[
\begin{aligned}
u_i&=\operatorname{unit}(\bar r_i^\top P_r),&
v_i&=\operatorname{unit}(\bar w_i^\top P_w),\\
k_i&=\operatorname{unit}(u_i\od v_i).&&
\end{aligned}
\]
The query and threshold use separate direct projections:
\[
\begin{aligned}
\hat q_a&=\operatorname{unit}(xP_a+b_a),\\
\tau_a&=-1+2\sigma(xu_a+\beta_a),&
\rho_i&=\hat q_a^\top k_i.
\end{aligned}
\]
For \texttt{linear\_angular} composition:
\[
c_i=
\operatorname{clip}\!\left(
\frac{\rho_i-\tau_a}{\max(1-\tau_a,10^{-4})},0,1\right)
\]
and padding rows are masked to zero. This is the single applicability
coefficient used in RW composition. Let \(M_a=\sum_i c_i\). The production
composed output before causal attention or residual addition is
\[
F_a(x)=
\sqrt{\frac{D}{L}}\,
\frac{\sum_i c_i\langle x,\bar r_i\rangle\bar w_i}
{\max(M_a,1)^{p_a}},
\]
with \(p_Q=p_K=0.5\), \(p_V=1.0\), and \(p_{\mathrm{RST}}=1.2\).
Operator rows are sharded over the model axis; local numerators and admission
mass are reduced across shards. Chunking changes the contraction schedule, not
the equation.

\section{Complete Model and Training Settings}
\label{app:settings}

\begin{table*}[!t]
\centering
\caption{Detailed settings from the registered training configurations and
restored checkpoint manifests.}
\scriptsize
\renewcommand{\arraystretch}{0.94}
\begin{tabular}{lrrr}
\toprule
Setting & \DAWN{} 400M & Dense 400M & \DAWN{} 1.3B\\
\midrule
Vocabulary & 30,522 & 30,522 & 30,522\\
Residual width & 2,304 & 1,280 & 4,096\\
Layers & 18 & 18 & 24\\
Heads / head width & 36 / 64 & 20 / 64 & 32 / 128\\
Maximum sequence length & 512 & 512 & 512\\
Feed-forward width & -- & 5,120 & --\\
Operation-address width & 256 & -- & 256\\
QK / V / RST pool & 4,600 / 14,500 / 29,300 & -- &
8,534 / 26,904 / 54,352\\
QK / V / RST chunks & 1 / 1 / 4 & -- & 1 / 2 / 8\\
QK / V / RST denominator power & 0.5 / 1.0 / 1.2 & -- &
0.5 / 1.0 / 1.2\\
Initial support-fraction targets & 0.05 / 0.03 / 0.01 & -- &
0.05 / 0.03 / 0.01\\
Training batch & 1,024 & 1,024 & 448\\
Mesh (data \(\times\) model) & \(16\times2\) & \(16\times2\) & \(16\times2\)\\
Learning rate & \(3\times10^{-4}\) & \(3\times10^{-4}\) & \(2\times10^{-4}\)\\
Warmup ratio & 0.08 & 0.08 & 0.06\\
Weight decay & 0.1 & 0.1 & 0.1\\
Tau learning-rate multiplier & 0.001 & -- & 0.001\\
Dropout / \texttt{router\_dropout} & 0 / 0 & 0 / -- & 0 / 0\\
Gradient checkpointing & yes & yes & yes\\
Seed & 1 & 1 & 1\\
\bottomrule
\end{tabular}
\end{table*}

The 400M \DAWN{} checkpoint is the primary performance and mechanistic model.
The dense 400M checkpoint is only a matched performance control. The 1.3B
\DAWN{} checkpoint is used for descriptive performance and one fail-closed IOI
localization replication. The still-incomplete dense 1.3B run is excluded.

\section{Benchmark Metrics and Protocol Details}
\label{app:benchmarks}

\smallhead{C4.}
All reported C4 rows use the same packed validation artifact and
\texttt{packed\_c4\_validation\_v1} protocol. The capped view contains 19,531
sequences. Exactly 19,520 sequences form 610 complete global evaluation
batches, producing 9,974,720 valid next-token targets; 11 incomplete-tail
sequences are excluded before scoring.

\smallhead{Zero-shot.}
The protocol is \texttt{mamba\_table3\_zero\_shot\_v1}, using stock
\texttt{lm-eval==0.4.2}, global batch 32, maximum length 512, length buckets
64/128/256/512, and zero few-shot examples. Evaluated split sizes are 5,153
for LAMBADA, 10,042 for HellaSwag, 1,838 for PIQA, 2,376 for ARC-Easy, 1,172
for ARC-Challenge, and 1,267 for WinoGrande. Task, dataset, tokenizer, and
software contracts match exactly between the two 400M models.

\section{IOI Localization and Intervention Details}
\label{app:ioi}

\subsection{Frozen discovery}

Discovery admits 119 paired-correct independent units and ranks 456,697 sites.
Split-half rank stability is 0.913978765. The frozen selection rule is the
smallest prefix satisfying cumulative absolute importance at least 0.70 and
split top-\(k\) overlap at least 0.95. The selected \(k=4096\) prefix has:
\begin{itemize}
    \item cumulative absolute importance 0.719248713;
    \item 3,993 shared split sites, overlap 0.974853516, and Jaccard
    0.950940700;
    \item Q/K/V/RST counts \(=1744/2350/2/0\);
    \item coverage of every layer from 0 through 17.
\end{itemize}
Validation and test do not recompute or alter this selection.

\begin{table*}[!t]
\centering
\caption{Control effects and frozen-minus-control separation. Each control
family contains 100 replicates. All paired permutation \(p\)-values are
0.0004997501 with 2,000 permutations.}
\label{tab:controls}
\scriptsize
\renewcommand{\arraystretch}{0.96}
\begin{tabular}{lrrrr}
\toprule
Control & Validation drop & Frozen minus control (95\% CI) &
Test drop & Frozen minus control (95\% CI)\\
\midrule
Equal-count & 0.070886 & 7.793902 [7.437851, 8.161572] &
0.080276 & 7.330057 [6.922408, 7.716879]\\
Same-layer & 0.039389 & 7.825399 [7.440629, 8.180288] &
0.041073 & 7.369260 [6.977447, 7.754539]\\
Activation-matched & 5.400490 & 2.464297 [2.076200, 2.832697] &
5.472579 & 1.937754 [1.564818, 2.319706]\\
Route-frequency & 0.219162 & 7.645626 [7.289512, 8.031625] &
0.242665 & 7.167668 [6.779836, 7.545537]\\
\bottomrule
\end{tabular}
\end{table*}

\subsection{Suppression, controls, and restoration}

For a correct candidate \(c\) and source candidate \(s\), the primary score is
\[
\mathcal{M}=\sum_{t\in c}\log p(t)-\sum_{t\in s}\log p(t).
\]
Suppression zeros the frozen sites' execution numerator at specified token
positions. It leaves admission mass and the production denominator unchanged.
Controls use the same 4,096-site count and exclude all circuit sites.
Activation matching conditions on layer, computation type, and the discovery mean
pre-scale operator-output-norm quantile. Restoration executes a distinct
intervention: after suppression, it adds back the exact selected same-example
numerator under the unchanged denominator.

On validation, restoration returns the margin to 3.462791, leaving residual
drop 0.002665. On held-out test, it returns the margin to 3.242905, leaving
residual drop 0.011815. The respective recovery fractions are 0.999661155
(95\% CI [0.998654544, 1.000584125]) and 0.998405645
(95\% CI [0.997323737, 0.999385266]).

\section{Negative Mechanistic Results}
\label{app:negative}

\smallhead{RAVEL.}
Discovery uses 142 independent units. Rank stability is 0.628762 for
Continent, 0.676243 for Country, and 0.660622 for Language, below the
preregistered requirement. No prefix passes the complete discovery rule, no
variable circuit is frozen, and validation, transfer, and test are not opened.

\smallhead{ARC.}
Discovery stability is 0.89386 and a 4,096-site Q/K-only circuit is frozen.
On 51 validation units, intact margin is 1.735534 and suppressed margin is
2.819139, giving a frozen margin drop of \(-1.083605\) (95\% CI
[\(-1.472180\), \(-0.690348\)]). The effect has the wrong sign and fails
control separation. Restoration recovery is 1.004083 (95\% CI
[0.994587, 1.013652]), but restoration alone does not pass the conjunction.
Held-out test remains closed.

\smallhead{1.3B IOI.}
The scale replication admits 125 independent units and passes rank stability
at 0.907187749. The largest allowed 4,096-site prefix reaches only 0.635144
cumulative absolute importance, below the 0.70 gate. No circuit is frozen;
validation and test remain closed.

\section{Compute Accounting}
\label{app:compute}

The applicability-support analysis evaluates 9,994,240 packed-C4 positions and
observes
address-stage margins before RW application. It stores aggregate statistics,
not per-token operator identities or vectors. Static accounting uses batch 1,
sequence length 512, and \(1\) MAC \(=2\) FLOPs. It counts major
matmul/einsum operations and omits softmax, normalization, gate elementwise
arithmetic, embedding lookup, residual additions, and communication
collectives.

\begin{table*}[!t]
\centering
\caption{Backbone static accounting. ``Exact applicability support'' retains full dense
address scoring and replaces only full-pool RW application with the measured
applicability support at positive applicability margin.}
\scriptsize
\renewcommand{\arraystretch}{0.96}
\begin{tabular}{lrrr}
\toprule
Component & Current dense \DAWN{} & Dense address + exact RW & Dense Transformer\\
\midrule
Address scoring + query/tau projections & 0.293742 & 0.293742 & --\\
RW application & 4.306187 & 0.383157 & --\\
Causal attention & 0.043487 & 0.043487 & 0.024159\\
Dense output / QKV / FFN maps & 0.097845 & 0.097845 & 0.362388\\
\midrule
Total TFLOPs & 4.741260 & 0.818230 & 0.386547\\
\bottomrule
\end{tabular}
\end{table*}

The estimate restricted to exact applicability support is mathematically
output-equivalent if the executor
can skip all and only zero-weight RW terms. It does not include indexed
address retrieval and does not demonstrate kernel parity. Instrumented
profiling time is excluded from latency claims because histogram collection is
part of the measured path.

\section{Reproducibility Identities}
\label{app:repro}

\begin{table*}[!t]
\centering
\caption{Primary immutable identities. Branch names are omitted because they
move; exact checkpoint steps, hashes, and artifact locations are the binding
identities.}
\scriptsize
\begin{tabularx}{\textwidth}{>{\raggedright\arraybackslash}p{0.13\textwidth}
>{\raggedright\arraybackslash}p{0.43\textwidth} Y}
\toprule
Evidence & Exact identity & Repository source\\
\midrule
\DAWN{} 400M & model \texttt{spatial-r1-v4.1.7.2}; step 76,293;
checkpoint identity \texttt{a7ce8afc...f91c17}; config hash
\texttt{08733ae4...5039b} &
\path{configs/train_config_v4172_400M_c4_40B_v4_64_ver1_den_qk0p5_v1p0_rst1p2.yaml}\\
Dense 400M & model \texttt{baseline-JAX}; step 76,293; config hash
\texttt{c3ce5d72...82fc25} &
\path{configs/train_config_baseline_tpu_400M_c4_40B_v4_64_tp2.yaml}\\
\DAWN{} 1.3B & model \texttt{spatial-r1-v4.1.7.2}; step 87,193;
checkpoint identity \texttt{2e38b41f...bf5f9}; config hash
\texttt{1e877263...4a5222} &
\path{configs/train_config_v4172_1p3B_c4_20B_v4_64_ver1_den_qk0p5_v1p0_rst1p2.yaml}\\
IOI circuit & parsed frozen-spec hash \texttt{6ec26618...418c};
site hash \texttt{573cfff8...0200}; selected-row hash
\texttt{1439f6f3...fe0c1} &
\path{configs/paper_ioi_frozen_circuit_v4172_400m.yaml}\\
Compute support & support hash \texttt{b1eac5f5...bdc2};
accounting hash \texttt{d7ad6115...984e} &
\path{assets/paper/v4172_400m_compute_support_step76293/}\\
\bottomrule
\end{tabularx}
\end{table*}

The machine-readable checkpoint registry is
\texttt{configs/paper\_checkpoint\_registry.yaml}. Exact evaluation commits,
software versions, data hashes, result roots, and complete artifact hashes are
recorded in the repository experiment ledger. Historical experiment branches
are evidence identities only; this source is maintained on the canonical main
branch.

\section*{Acknowledgments}

The author thanks Anthropic Claude, Claude Code, OpenAI ChatGPT, and OpenAI
Codex for assistance throughout this research. This work was developed with
extensive use of AI systems for literature exploration, conceptual exploration,
cross-domain framing, abstraction and terminology refinement, implementation
support, debugging, and manuscript preparation. In particular, these systems
were used as interactive assistants for refining the learned state-transition
formulation and the exposition of \DAWN{}. The research direction,
architectural decisions, experimental design, experiments, analyses,
interpretations, and final claims remain the responsibility of the author.

This research was supported by the Google TPU Research Cloud (TRC) program,
which provided computational resources for training and analysis.

\begingroup
\footnotesize
\setlength{\bibsep}{0pt}
\bibliography{references}
\bibliographystyle{icml2026}
\endgroup

\end{document}
